Gradient direction during ERM training encodes which samples conflict with the dominant batch learning signal --- a property theoretically well-suited for identifying spurious-minority group members without group annotations.
We test this directly: we compute per-sample last-layer gradient cosine similarity with the within-batch mean gradient and measure its ROC-AUC as a spurious-minority detector on Waterbirds and CelebA across training epochs $\{1, 5, 10, 25, 50\}$.
The signal fails decisively: alignment ROC-AUC reaches only 0.35 on Waterbirds and 0.63 on CelebA at best, far below per-sample loss (0.77--0.97, across datasets and epochs), and is \emph{inverted} on Waterbirds at epoch~1 (ROC-AUC = 0.15), where spurious-minority samples appear more aligned with the batch mean than majority samples.
We attribute this to \textbf{batch contamination}: high-loss minority samples disproportionately pull the within-batch mean gradient toward their own direction, paradoxically increasing their measured cosine similarity with the mean.
This failure mode is prevalence-dependent, explaining why the inversion is weaker on CelebA (${\sim}0.8\%$ minority) than Waterbirds (${\sim}5\%$).
We derive an empirically grounded design principle: two-pass global mean gradient reference (not within-batch mean) is expected to be necessary to avoid contamination and recover the directional minority signal.
Our results constitute the first direct empirical characterization of gradient alignment ROC-AUC as a spurious-minority detector on these benchmarks, and provide a concrete corrective direction for gradient-based annotation-free debiasing methods.
